\documentclass[11pt]{amsart}
\usepackage{amsmath,amssymb,amsthm}
\usepackage[hidelinks]{hyperref}

\newtheorem{theorem}{Theorem}
\newtheorem{lemma}{Lemma}
\newtheorem{corollary}{Corollary}
\theoremstyle{definition}
\newtheorem{problem}{Problem}
\newtheorem{remark}{Remark}

\newcommand{\ellone}{\ell_1}
\newcommand{\spn}{\operatorname{span}}
\newcommand{\denspan}{\overline{\operatorname{span}}}
\newcommand{\supp}{\operatorname{supp}}

\title{Closed sublattices of \(\ell_1(\Gamma)\) in the norming-M-basis heredity problem}
\author{FA Banach run packet}
\date{June 16, 2026}

\begin{document}
\maketitle

\section*{Status}

\textbf{Partial result.}
Closed sublattices of \(\ell_1(\Gamma)\) have \(1\)-norming Markushevich
bases. Consequently, closed subspaces of \(\ell_1(\Gamma)\) that are finite
codimensional perturbations of closed sublattices have norming M-bases.

\section*{Source Problem}

Hajek and Russo ask the following heredity question for norming
Markushevich bases \cite[Problem 6.1]{HR}. The problem appears in
Section 6 of arXiv:2402.18442.

\begin{problem}[Hajek--Russo]
Let \(X\) be a Banach space with norming M-basis and let \(Z\) be a
subspace of \(X\). Must \(Z\) admit a norming M-basis as well?
\end{problem}

The source identifies subspaces of \(\ell_1(\Gamma)\), \(\Gamma\)
uncountable, as the simplest possible counterexample territory and notes
that it is not even known whether all such subspaces are Plichko. It also
records that hyperplanes and closed sublattices of \(\ell_1(\Gamma)\) are
known to be \(1\)-Plichko \cite[Section 6]{HR}.

\section*{Idea of the Proof}

A closed sublattice \(L\) of \(\ell_1(\Gamma)\) is much more rigid than a
general closed subspace. Its norm is additive on positive vectors, so \(L\)
is an AL-space. Kakutani's representation theorem writes \(L\) as an
\(L_1(\mu)\) space. The coordinate functionals of the ambient
\(\ell_1(\Gamma)\), restricted to \(L\), are lattice homomorphisms and
separate points of \(L\). In an \(L_1(\mu)\) space, non-zero lattice
homomorphisms live on atoms. Since the ambient coordinates separate points,
there can be no nonatomic part. Thus \(L\) is an atomic \(L_1\) space,
hence a weighted \(\ell_1(A)\). Its atom basis is the required
\(1\)-norming M-basis.

\section*{The Result}

\begin{lemma}\label{lem:sublattice-atomic}
Let \(L\) be a closed sublattice of \(\ell_1(\Gamma)\). Then \(L\) is
lattice isometric to \(\ell_1(A,w)\) for some set \(A\) and weights
\(w_a>0\). In particular, \(L\) is linearly isometric to \(\ell_1(A)\).
\end{lemma}

\begin{proof}
The lattice operations on \(L\) are inherited coordinatewise from
\(\ell_1(\Gamma)\). If \(u,v\in L_+\), then
\[
  \|u+v\|_1=\sum_{\gamma\in\Gamma}(u_\gamma+v_\gamma)
           =\|u\|_1+\|v\|_1.
\]
Thus \(L\) is an AL-space. By Kakutani's representation theorem for
AL-spaces, \(L\) is lattice isometric to \(L_1(\mu)\) for some measure
space \((\Omega,\Sigma,\mu)\); see, for instance, the standard formulation
quoted in \cite[Theorem 4.27]{AB}.

For each \(\gamma\in\Gamma\), let
\[
  h_\gamma=\delta_\gamma|_L .
\]
This functional is either zero or a lattice homomorphism: for example,
\[
  h_\gamma(x\vee y)=\max\{h_\gamma(x),h_\gamma(y)\}
\]
and similarly for \(x\wedge y\), since the operations are coordinatewise.
The family \(\{h_\gamma:\gamma\in\Gamma\}\) separates points of \(L\).

Transport these functionals through the lattice isometry
\(L\simeq L_1(\mu)\). We use the elementary description of lattice
homomorphisms on \(L_1\): a non-zero continuous lattice homomorphism on
\(L_1(\mu)\) is supported on a single atom of \(\mu\). Indeed, such a
homomorphism is a positive functional, hence is represented by a positive
measure absolutely continuous with respect to \(\mu\). If it gave positive
mass to two disjoint measurable pieces, the identity
\(\varphi(f\wedge g)=\min\{\varphi(f),\varphi(g)\}\) applied to positive
disjoint \(f,g\) would be violated. Therefore its representing measure is
concentrated on one atom.

Since the transported coordinate homomorphisms separate points of
\(L_1(\mu)\), every non-zero \(f\in L_1(\mu)\) must be detected by an atom.
Hence the nonatomic part of \(\mu\) is null. Therefore \(L_1(\mu)\) is an
atomic \(L_1\) space, i.e. lattice isometric to a weighted
\(\ell_1(A,w)\). Rescaling the atom vectors gives a linear isometry with
\(\ell_1(A)\).
\end{proof}

\begin{theorem}\label{thm:closed-sublattice}
Every closed sublattice \(L\) of \(\ell_1(\Gamma)\) admits a
\(1\)-norming Markushevich basis.
\end{theorem}

\begin{proof}
By Lemma \ref{lem:sublattice-atomic}, \(L\) is linearly isometric to
\(\ell_1(A)\). The canonical unit vectors \((e_a)_{a\in A}\), together
with the coordinate functionals \((e_a^*)_{a\in A}\), form a fundamental
biorthogonal system. The linear span of the coordinates is \(1\)-norming:
for every \(x\in \ell_1(A)\),
\[
  \|x\|_1
  =\sup\left\{\left|\sum_{a\in F}\varepsilon_a x_a\right|:
       F\subset A \text{ finite},\ |\varepsilon_a|\le 1\right\}.
\]
Thus the canonical system is a \(1\)-norming M-basis, and transporting it
through the isometry gives a \(1\)-norming M-basis of \(L\).
\end{proof}

\begin{corollary}\label{cor:finite-commensurable}
Let \(Z\) be a closed subspace of \(\ell_1(\Gamma)\). Suppose there is a
closed sublattice \(L\subseteq \ell_1(\Gamma)\) such that, for
\(Y=Z+L\),
\[
  \dim(Y/L)<\infty
  \quad\text{and}\quad
  \dim(Y/Z)<\infty .
\]
Then \(Z\) admits a norming M-basis.
\end{corollary}

\begin{proof}
Theorem \ref{thm:closed-sublattice} gives a norming M-basis for \(L\).
The finite-codimensional invariance theorem promoted in the companion run
packet \cite{FiniteCodimPacket} says that a Banach space has a norming
M-basis if and only if each of its closed finite-codimensional subspaces
does. Since \(Y/L\) is finite-dimensional, \(L\) is finite-codimensional
in \(Y\), so \(Y\) has a norming M-basis. Since \(Y/Z\) is
finite-dimensional, \(Z\) is finite-codimensional in \(Y\), and the same
invariance theorem gives a norming M-basis for \(Z\).
\end{proof}

\section*{Limitations}

The proof uses that \(L\) is a sublattice. Arbitrary closed subspaces of
\(\ell_1(\Gamma)\) do not inherit the coordinate lattice operations, and
the argument gives no atomic representation for them. Thus the main
subspace problem for \(\ell_1(\Gamma)\), and hence Problem 6.1 in general,
remains open.

\section*{Novelty and Literature Check}

The source paper's arXiv page lists version 3, revised September 30, 2025,
and still presents Problem 6.1 as open \cite{HRarxiv}. Web searches for
the exact heredity problem, norming M-bases in subspaces, Kalenda's
\(\ell_1(\Gamma)\) Plichko question, and closed sublattices of
\(\ell_1(\Gamma)\) found no later full solution and no exact
closed-sublattice norming-M-basis statement. The packet should be read as
a standard Banach-lattice observation applied to the source problem, not
as a claim that the arbitrary \(\ell_1(\Gamma)\) subspace problem is
settled.

\begin{thebibliography}{9}

\bibitem{AB}
C.~D. Aliprantis and O.~Burkinshaw,
\emph{Positive Operators},
Springer, 2006.

\bibitem{HR}
P.~Hajek and T.~Russo,
\emph{Norming Markushevich bases: recent results and open problems},
arXiv:2402.18442.

\bibitem{HRarxiv}
P.~Hajek and T.~Russo,
arXiv page for \emph{Norming Markushevich bases: recent results and open
problems}, \url{https://arxiv.org/abs/2402.18442}.

\bibitem{FiniteCodimPacket}
FA Banach run packet,
\emph{Finite-codimensional heredity for norming Markushevich bases},
\texttt{runs/fa\_banach\_001/solutions/partial/2402.18442\_finite\_codim\_heredity\_norming\_mbasis}.

\end{thebibliography}

\end{document}
